% GENERATED FILE. Edit canonical lessons, then run npm run generate:books.
% canonical-source-hash: fnv1a64:4c9080726dec612a
% canonical-lessons: PT-C67-barato, PT-C67-alto, PT-C67-baixo, PT-C67-cheio, PT-C67-vazio

\chapter{Cheap, Tall, Low, Full, Empty}
\label{ch:pt-cheap-tall-low-full-empty}
\hlchaptermodality{67}

\begin{chapteropening}
\textbf{By the end of this chapter:} \emph{I can say cheap, tall, low, full and empty.}

Complete the last lesson of chapter 67: I can say cheap, tall, low, full and empty.
\end{chapteropening}

\section[barato]{barato — cheap}

\label{lesson:PT-C67-barato}



\begin{quote}
Before the new one: say the Portuguese for difficult, then the Portuguese for expensive; dear.
\end{quote}

\subsection*{You'll want to know: barato}
\textbf{barato} — \textquotedblleft{}cheap\textquotedblright{}. The opposite of \textbf{caro}.

\begin{cousinweb}
\textbf{barato} comes from an old verb for trading and selling low, whose own origin is uncertain. Old French \emph{barater}, behind English \emph{barter}, belongs to the same family.
\end{cousinweb}

\begin{grammarlens}[title={What you've built}]
One more word for this chapter.
\end{grammarlens}

\subsection*{Your turn}
\begin{itemize}
\raggedright
  \item \emph{Say it:} \emph{barato}
  \item \emph{Say it:} \emph{barato}, once more
  \item \emph{Say it:} say \emph{barato}
  \item \emph{Recall:} say the Portuguese for difficult, then the Portuguese for expensive; dear, then say \emph{barato} again
\end{itemize}

\begin{tcolorbox}[breakable,colback=teal!4,colframe=teal!35!black,title={Before you move on}]
What does barato mean? (\textquotedblleft{}Cheap\textquotedblright{}.) Say it once more.
\end{tcolorbox}
\section[alto]{alto — tall, high}

\label{lesson:PT-C67-alto}



\begin{quote}
Before the new one: say the Portuguese for expensive; dear, then the Portuguese for cheap.
\end{quote}

\subsection*{You'll want to know: alto}
\textbf{alto} — \textquotedblleft{}tall, high\textquotedblright{}.

\begin{cousinweb}
\textbf{alto} $\leftarrow$ Latin \emph{altus}, \textquotedblleft{}high, deep\textquotedblright{} (EtymDB 2.1) — inherited.
\end{cousinweb}

\begin{grammarlens}[title={What you've built}]
One more word for this chapter.
\end{grammarlens}

\subsection*{Your turn}
\begin{itemize}
\raggedright
  \item \emph{Say it:} \emph{alto}
  \item \emph{Say it:} \emph{alto}, once more
  \item \emph{Say it:} say \emph{alto}
  \item \emph{Recall:} say the Portuguese for expensive; dear, then the Portuguese for cheap, then say \emph{alto} again
\end{itemize}

\begin{tcolorbox}[breakable,colback=teal!4,colframe=teal!35!black,title={Before you move on}]
What does alto mean? (\textquotedblleft{}Tall, high\textquotedblright{}.) Say it once more.
\end{tcolorbox}
\section[baixo]{baixo — short, low}

\label{lesson:PT-C67-baixo}



\begin{quote}
Before the new one: say the Portuguese for cheap, then the Portuguese for tall.
\end{quote}

\subsection*{You'll want to know: baixo}
\textbf{baixo} — \textquotedblleft{}short, low\textquotedblright{}.

\begin{cousinweb}
\textbf{baixo} $\leftarrow$ Late Latin \emph{bassus}, \textquotedblleft{}low\textquotedblright{} (EtymDB 2.1) — inherited.
\end{cousinweb}

\begin{grammarlens}[title={What you've built}]
One more word for this chapter.
\end{grammarlens}

\subsection*{Your turn}
\begin{itemize}
\raggedright
  \item \emph{Say it:} \emph{baixo}
  \item \emph{Say it:} \emph{baixo}, once more
  \item \emph{Say it:} say \emph{baixo}
  \item \emph{Recall:} say the Portuguese for cheap, then the Portuguese for tall, then say \emph{baixo} again
\end{itemize}

\begin{tcolorbox}[breakable,colback=teal!4,colframe=teal!35!black,title={Before you move on}]
What does baixo mean? (\textquotedblleft{}Short, low\textquotedblright{}.) Say it once more.
\end{tcolorbox}
\section[cheio]{cheio — full}

\label{lesson:PT-C67-cheio}



\begin{quote}
Before the new one: say the Portuguese for tall, then the Portuguese for short.
\end{quote}

\subsection*{You'll want to know: cheio}
\textbf{cheio} — \textquotedblleft{}full\textquotedblright{}.

\begin{cousinweb}
\textbf{cheio} $\leftarrow$ Latin \emph{plenus}, \textquotedblleft{}full\textquotedblright{} (EtymDB 2.1) — inherited. Latin \emph{pl-} became \emph{ch-}, and the \emph{-n-} fell out.
\end{cousinweb}

\begin{grammarlens}[title={What you've built}]
One more word for this chapter.
\end{grammarlens}

\subsection*{Your turn}
\begin{itemize}
\raggedright
  \item \emph{Say it:} \emph{cheio}
  \item \emph{Say it:} \emph{cheio}, once more
  \item \emph{Say it:} say \emph{cheio}
  \item \emph{Recall:} say the Portuguese for tall, then the Portuguese for short, then say \emph{cheio} again
\end{itemize}

\begin{tcolorbox}[breakable,colback=teal!4,colframe=teal!35!black,title={Before you move on}]
What does cheio mean? (\textquotedblleft{}Full\textquotedblright{}.) Say it once more.
\end{tcolorbox}
\section[vazio]{vazio — empty}

\label{lesson:PT-C67-vazio}



\begin{quote}
Before the new one: say the Portuguese for short, then the Portuguese for full.
\end{quote}

\subsection*{You'll want to know: vazio}
\textbf{vazio} — \textquotedblleft{}empty\textquotedblright{}.

\begin{cousinweb}
\textbf{vazio} $\leftarrow$ Latin \emph{vacivus}, \textquotedblleft{}empty\textquotedblright{} (EtymDB 2.1), from \emph{vacare}, \textquotedblleft{}to be empty\textquotedblright{} — inherited.
\end{cousinweb}

\begin{grammarlens}[title={What you've built}]
One more word for this chapter.
\end{grammarlens}

\subsection*{Your turn}
\begin{itemize}
\raggedright
  \item \emph{Say it:} \emph{vazio}
  \item \emph{Say it:} \emph{vazio}, once more
  \item \emph{Say it:} say \emph{vazio}
  \item \emph{Recall:} say the Portuguese for short, then the Portuguese for full, then say \emph{vazio} again
\end{itemize}

\begin{tcolorbox}[breakable,colback=teal!4,colframe=teal!35!black,title={Before you move on}]
What does vazio mean? (\textquotedblleft{}Empty\textquotedblright{}.) Say it once more.
\end{tcolorbox}
